\documentclass[../../../include/open-logic-section]{subfiles}

\begin{document}

\olfileid{sth}{ordinals}{replacement}
\olsection{الاستبدال}

كان وجه إدخال الأعداد الترتيبية في~\olref[sth][ordinals][vn]{sec}
أنها تصلح أنماطًا للترتيب، أي ممثلات قياسية للترتيبات الجيدة. ولا يتم
هذا إلا بإثبات أن \emph{كل ترتيب جيد يتماثل مع عدد ترتيبي ما}.
فعندئذ نعرّف~$\ordtype{{\OLId{latin-looped}{A}}, <}$ بأنه العدد الترتيبي~${\OLId{greek-variable}{alpha}}$ الذي
يحقق~$\tuple{{\OLId{latin-looped}{A}}, <} \isomorphic {\OLId{greek-variable}{alpha}}$.

غير أن المسلّمات المتقدمة وحدها \emph{لا تكفي} لإثبات هذه النتيجة.
وسيأتي بيان العجز في~\olref[replacement][strength]{sec}؛ وحسبنا الآن
أن نعرف أن البرهان لا يتم بها. فلا بد من أصل جديد، وهو الآتي.

\begin{axiom}[مخطط الاستبدال]
إذا أُعطيت صيغة~$\phi({\OLId{latin-ordinary}{x}}, {\OLId{latin-ordinary}{y}})$، عُدّ القول الآتي مسلّمة:
\begin{quote}
لكل~${\OLId{latin-looped}{A}}$، إذا كان~$(\forall {\OLId{latin-ordinary}{x}} \in {\OLId{latin-looped}{A}})\lexists![{\OLId{latin-ordinary}{y}}][\phi({\OLId{latin-ordinary}{x}},{\OLId{latin-ordinary}{y}})]$،
وجدت المجموعة~$\Setabs{{\OLId{latin-ordinary}{y}}}{(\exists {\OLId{latin-ordinary}{x}} \in {\OLId{latin-looped}{A}})\phi({\OLId{latin-ordinary}{x}},{\OLId{latin-ordinary}{y}})}$.
\end{quote}
\end{axiom}
\noindent
وهذا مخطط كالفصل، لا مسلّمة مفردة: فإن الصيغ~$\phi$ غير متناهية
العدد، ولكل منها مسلّمة. وله صياغة أخرى مكافئة، هي المعتادة:

\begin{defish}
لكل صيغة~$\phi({\OLId{latin-ordinary}{x}},{\OLId{latin-ordinary}{y}})$ لا يرد فيها «${\OLId{latin-looped}{B}}$»، تُؤخذ الصيغة الآتية مسلّمة:
\[
\forall {\OLId{latin-looped}{A}}[(\forall {\OLId{latin-ordinary}{x}} \in {\OLId{latin-looped}{A}})\lexists![{\OLId{latin-ordinary}{y}}][\phi({\OLId{latin-ordinary}{x}},{\OLId{latin-ordinary}{y}})] \lif
\exists {\OLId{latin-looped}{B}}\forall {\OLId{latin-ordinary}{y}} ({\OLId{latin-ordinary}{y}} \in {\OLId{latin-looped}{B}} \liff (\exists {\OLId{latin-ordinary}{x}} \in {\OLId{latin-looped}{A}})\phi({\OLId{latin-ordinary}{x}},{\OLId{latin-ordinary}{y}}))]
\]
\end{defish}

وقد يثقل فهم هذه الصيغة لأول وهلة لتشابك أجزائها. أما النتيجة
القريبة التالية فلعلها \emph{أبين} في الكشف عن المعنى الحدسي المقصود.
\footnote{الحد تعبير يعيّن شيئًا واحدًا لا غير كلما عُيّنت قيم
متغيراته الحرة. فـ«$\emptyset$» يعيّن المجموعة الخالية؛ و«$\{{\OLId{latin-ordinary}{x}}\}$»
يعيّن أحادية~${\OLId{latin-ordinary}{x}}$، مهما يكن~${\OLId{latin-ordinary}{x}}$؛ و«${\OLId{latin-ordinary}{x}} \cup {\OLId{latin-ordinary}{y}}$» يعيّن اتحاد~${\OLId{latin-ordinary}{x}}$
و${\OLId{latin-ordinary}{y}}$، أيًا كانا.}

\begin{cor}
لكل حد~$\tau({\OLId{latin-ordinary}{x}})$ ومجموعة~${\OLId{latin-looped}{A}}$ توجد المجموعة:
\[
\Setabs{\tau({\OLId{latin-ordinary}{x}})}{{\OLId{latin-ordinary}{x}} \in {\OLId{latin-looped}{A}}} = \Setabs{{\OLId{latin-ordinary}{y}}}{(\exists {\OLId{latin-ordinary}{x}} \in {\OLId{latin-looped}{A}}){\OLId{latin-ordinary}{y}} = \tau({\OLId{latin-ordinary}{x}})}.
\]
\end{cor}

\begin{proof}
بما أن~$\tau$ \emph{حد}، فإن~$\forall {\OLId{latin-ordinary}{x}} \lexists![{\OLId{latin-ordinary}{y}}][\tau({\OLId{latin-ordinary}{x}}) = {\OLId{latin-ordinary}{y}}]$،
ومن باب أولى~$(\forall {\OLId{latin-ordinary}{x}} \in {\OLId{latin-looped}{A}})\lexists![{\OLId{latin-ordinary}{y}}][\tau({\OLId{latin-ordinary}{x}}) = {\OLId{latin-ordinary}{y}}]$.
فيقضي الاستبدال بوجود~$\Setabs{{\OLId{latin-ordinary}{y}}}{(\exists {\OLId{latin-ordinary}{x}} \in {\OLId{latin-looped}{A}})\tau({\OLId{latin-ordinary}{x}}) = {\OLId{latin-ordinary}{y}}}$.
\end{proof}
\noindent
ومن هنا تظهر مناسبة اسم «الاستبدال»: نبدأ بمجموعة~${\OLId{latin-looped}{A}}$، ثم ننشئ
المجموعة~$\Setabs{\tau({\OLId{latin-ordinary}{x}})}{{\OLId{latin-ordinary}{x}} \in {\OLId{latin-looped}{A}}}$ بأن نستبدل بكل عنصر من~${\OLId{latin-looped}{A}}$
صورته تحت~$\tau$. وقياسًا على ترميز صورة المجموعة تحت الدالة،
يجوز أن نرمز بـ$\funimage{\tau}{{\OLId{latin-looped}{A}}}$ إلى~$\Setabs{\tau({\OLId{latin-ordinary}{x}})}{{\OLId{latin-ordinary}{x}} \in {\OLId{latin-looped}{A}}}$.

ولا يغيبنّ الفرق الحاسم: إن~$\tau$ \emph{حد}، وليس من شرطه أن
يسمي \emph{دالة} بالمعنى المحدد في~\olref[sfr][fun][rel]{sec}،
أي مجموعة من الأزواج المرتبة. فلو كانت~${\OLId{latin-ordinary}{f}}$ دالة بهذا المعنى، لكانت
$\funimage{{\OLId{latin-ordinary}{f}}}{{\OLId{latin-looped}{A}}} = \Setabs{{\OLId{latin-ordinary}{f}}({\OLId{latin-ordinary}{x}})}{{\OLId{latin-ordinary}{x}} \in {\OLId{latin-looped}{A}}}$ مجرد مجموعة جزئية من
$\ran{{\OLId{latin-ordinary}{f}}}$، ووجودها مضمون بمسلّمات~$\Zminus$ وحدها.
\footnote{يكفي أن نأخذ
$\Setabs{{\OLId{latin-ordinary}{y}} \in \bigcup \bigcup {\OLId{latin-ordinary}{f}}}{(\exists {\OLId{latin-ordinary}{x}} \in {\OLId{latin-looped}{A}}){\OLId{latin-ordinary}{y}} = {\OLId{latin-ordinary}{f}}({\OLId{latin-ordinary}{x}})}$.}
أما الاستبدال فيزيد مسلّماتنا \emph{قوة}، كما سيظهر في
\olref[sth][replacement][]{chap}.

\end{document}
